\subsection{Sensitivity of focus displacement}\label{supp:fem-sensitivities}
The spatial ordering in the article persists across montage, labeling, tissue,
and focus-fraction choices. Table~\ref{tab:fem-sensitivities} gives the ranges;
orientation changes absolute contrast more than the atlas label in this subject.
For the reference projection and label, the top-5--50\% focus moves
$2.2$--$4.0$~mm anteriorly for the envelope and $2.4$--$8.2$~mm posteriorly
for mean-square polarization.

The sweep shows two kinds of steering (Fig.~\ref{fig:fem}e,f). The envelope ratio peaks
near balance ($1.24$ at a posterior share of $0.525$) and falls to $0.89$--$0.94$ at
the extremes. Its focus moves only while the two carriers are comparable inside the
hippocampus. At a posterior share of $0.6$ or more, the anterior pair is the weaker
carrier at more than 96\% of hippocampal nodes; at $0.4$ or less, the posterior pair is. The
envelope map then equals twice the weaker pair's polarization over most of the ROI,
and its measured focus changes little with further imbalance. The focus of $\mathcal I$ moves
continuously instead, from $1.5$~mm anterior at a posterior share of $0.1$ to $6.3$~mm
posterior at $0.9$, while its ratio stays between $0.79$ and $0.94$ ($0.89$--$0.97$
degree-matched). The quadratic excursion keeps one ratio and one focus across the
whole sweep.

\begin{table}[htbp]\centering\small
\caption{Focus displacement from 1:1 to 1:3. Envelope shifts are anterior;
mean-square shifts are posterior. The quadratic beat shift is zero in every case.
Contrasts quoted in the last column are raw 1:1 envelope/quadratic contrasts.}
\label{tab:fem-sensitivities}
\begin{tabular}{@{}lrrl@{}}
\toprule
Sensitivity & Envelope (mm) & Mean square (mm) & Contrasts\\
\midrule
Atlas label & $2.0$--$3.8$ & $3.2$--$5.9$ & $1.17$ / $1.04$\\
Radial direction & $1.9$--$2.8$ & $2.0$--$3.7$ & $1.62$ / $1.97$\\
All 32 combinations & $0.7$--$8.1$ & $0.5$--$10.6$ & \\
\bottomrule
\end{tabular}
\end{table}


The gray-matter-only comparison uses 1,174 FreeSurfer hippocampal nodes
against the unchanged gray-matter reference. Table~\ref{tab:fem-gm-contrasts}
gives both contrast definitions for all four maps. Degree matching takes
the square root of a map before computing its statistics. At 1:1 only the
envelope p95 contrast exceeds one; at 1:3 none does. No raw or degree-matched
mean contrast reaches one in either tissue definition.

\begin{table}[htbp]\centering\small
\caption{Gray-matter-only hippocampal contrasts, 1:1 / 1:3. Each denominator
is the reference p95 of the same raw or square-root map.}
\label{tab:fem-gm-contrasts}
\setlength{\tabcolsep}{4pt}
\begin{tabular}{@{}lrrrr@{}}
\toprule
& \multicolumn{2}{c}{Raw map} & \multicolumn{2}{c}{Square-root map}\\
Map & p95 contrast & Mean contrast & p95 contrast & Mean contrast\\
\midrule
Envelope & $1.086/0.839$ & $0.757/0.547$ & -- & --\\
Quadratic beat & $0.935/0.935$ & $0.508/0.508$ & $0.967/0.967$ & $0.696/0.696$\\
Isotropic beat & $0.943/0.943$ & $0.512/0.512$ & $0.971/0.971$ & $0.698/0.698$\\
Mean square & $0.715/0.765$ & $0.396/0.462$ & $0.845/0.874$ & $0.615/0.665$\\
\bottomrule
\end{tabular}
\end{table}
